\section{高階の線形漸化式}\label{節：高階の線形漸化式}
\subsection{過去を状態にまとめる}

次の項に過去三項が必要なら,\,その三項を状態として持つ.
例えば
\[
 a_{n+3}=6a_{n+2}-11a_{n+1}+6a_n
\]
は
\[
 \begin{pmatrix}a_{n+3}\\a_{n+2}\\a_{n+1}\end{pmatrix}
 =
 \begin{pmatrix}6&-11&6\\1&0&0\\0&1&0\end{pmatrix}
 \begin{pmatrix}a_{n+2}\\a_{n+1}\\a_n\end{pmatrix}.
\]
第一行で新しい項を作り,\,下の行で記録を一つずつ移す.
二成分の場合から増えたのは,\,記憶する過去の長さである.

一般に定数係数の$k$階線形漸化式
\[
 a_{n+k}=\sum_{j=0}^{k-1}c_j a_{n+j}
\]
に対して
\[
 \boldsymbol v_n=
 \begin{pmatrix}a_{n+k-1}\\a_{n+k-2}\\\vdots\\a_n\end{pmatrix}
\]
と置くと,$\boldsymbol v_{n+1}=C\boldsymbol v_n$である.
更新行列は
\[
 C=
 \begin{pmatrix}
 c_{k-1}&c_{k-2}&\cdots&c_1&c_0\\
 1&0&\cdots&0&0\\
 0&1&\cdots&0&0\\
 \vdots&&\ddots&&\vdots\\
 0&0&\cdots&1&0
 \end{pmatrix}.
\]
成分数が増えても,\,ベクトルの和や行列の積は各行ごとに同じ方法で定める.
この形の行列を\textbf{コンパニオン行列}という.

\subsection{高階の特性方程式}

固有ベクトルの一番下の成分を1にそろえると,\,下の行から順に
\[
 \boldsymbol u(\lambda)=
 \begin{pmatrix}\lambda^{k-1}\\\lambda^{k-2}\\\vdots\\1\end{pmatrix}
\]
という形になる.
一番下が零なら上も順に零となるため,\,このそろえ方ができる.
第一行の条件から
\[
 \lambda^k=\sum_{j=0}^{k-1}c_j\lambda^j
\]
を得る. 従って特性多項式は
\[
 \chi(X)=X^k-\sum_{j=0}^{k-1}c_jX^j.
\]
高階でも,\,特性方程式$\chi(\lambda)=0$は一定倍率の方向を探している.
高次方程式の根を簡単な式で求められるかどうかと,\,この構造が成り立つかどうかは別の問題である.

\noindent\textbf{【例題】}\par\nobreak
$a_1=1,\,a_2=2,\,a_3=6$とし,
\[
 a_{n+3}=6a_{n+2}-11a_{n+1}+6a_n
\]
の一般項を求めよ.

\noindent\textbf{【方針】}\par\nobreak
特性多項式を因数分解し,\,三つの等比的な成分に分けます.

\noindent\textbf{【解答】}\par\nobreak
\[
 a_n=1-2^{n-1}+3^{n-1}.
\]
\noindent\textbf{【解法】}\par\nobreak
\[
 X^3-6X^2+11X-6=(X-1)(X-2)(X-3).
\]
従って$a_n=A+B2^{n-1}+D3^{n-1}$と置く.
初めの三項から
\[
 A+B+D=1,\quad A+2B+3D=2,
\]
\[
 A+4B+9D=6.
\]
これを解くと$A=1,\,B=-1,\,D=1$となる.
得た式は初期三項と更新を満たすので,\,全項に一致する.

\subsection{初期条件と解の自由度}\label{節：解の自由度}

$k$階の更新では,\,初めの$k$項を自由に選ぶと,\,その後は順に一意に決まる.
従って解の数列全体は,\,初期値の組$\mathbb C^k$と一対一に対応する.
二つの解の和と定数倍も解であるため,\,初期値だけでなく解全体にも線形性がある.
ここで「解全体は$k$次元」というのは,\,\textbf{独立に選べる初期の数が$k$個ある}という意味である.

特性根$\lambda_1,\ldots,\lambda_k$が相異なり非零なら,
\[
 a_n=\sum_{j=1}^{k}A_j\lambda_j^{n-1}
\]
で全解を表せる.
なぜこれで十分なのかを確かめよう.
各冪は更新を満たす. また,\,異なる冪の数列は一次独立である.

その独立性は,\,項を一つ進める操作$E$を使うと示しやすい.
$(Ea)_n=a_{n+1}$と定めると,
\[
 (E-\mu)\lambda^{n-1}
 =(\lambda-\mu)\lambda^{n-1}.
\]
もし$\sum_j A_j\lambda_j^{n-1}$が零数列なら,\,固定した$j$以外の根について
$E-\lambda_\ell$をすべて作用させる.
$j$以外の項が消え,
\[
 A_j\prod_{\ell\ne j}(\lambda_j-\lambda_\ell)
 \lambda_j^{n-1}=0
\]
となる. $n=1$を見れば$A_j=0$である.
根が異なるので積は零でない.

従って$k$本は独立で,\,自由度$k$の解全体の基底になる.
初期値を合わせる係数も一意に決まる.
一般項に必要な任意定数の個数は,\,記憶すべき初期条件の個数と対応する.

\subsection{重根と高階差}

非零の根$r$が重なると,\,冪$r^{n-1}$だけでは自由度が不足する.
そこで,\,多項式$P(n)$を掛けた
\[
 a_n=r^{n-1}P(n)
\]
を調べる. $\Delta P(n)=P(n+1)-P(n)$を\textbf{前進差分}と書くと,
\[
 (E-r)a_n=r^n\Delta P(n).
\]
もう一度作用させれば$r^{n+1}\Delta^2P(n)$になる.
一般に
\[
 (E-r)^m a_n=r^{n+m-1}\Delta^mP(n).
\]
差分を一回取ると非定数多項式の次数が一つ下がるので,\,次数$m-1$以下なら$\Delta^mP=0$である.
従って,$r$が$m$重根なら
\[
 r^{n-1},\ (n-1)r^{n-1},\ \ldots,\
 (n-1)^{m-1}r^{n-1}
\]
が解になる.

すべての根が非零で,\,異なる根を$r_1,\ldots,r_h$,
その重複度を$m_1,\ldots,m_h$とすれば,
\[
 a_n=\sum_{j=1}^{h}r_j^{n-1}P_j(n),
 \qquad \deg P_j<m_j
\]
で全解を表せる.
ここで$\sum_jm_j=k$であり,\,係数の総数も$k$である.
これらの独立性も差分で示せる.
ある根$r_j$に注目し,\,他の根の因子$(E-r_\ell)^{m_\ell}$を作用させると他の成分は消える.
一方,$r_j$の多項式成分では,\,異なる根の因子は最高次係数を非零倍するため次数を保つ.
残った成分が零数列ならその多項式は零であり,\,各根について同じことがいえる.
独立な解が$k$本あるので全解を尽くす.

例えば
\[
 a_{n+3}-3a_{n+2}+3a_{n+1}-a_n=0
\]
では,\,特性多項式は$(X-1)^3$である.
三階差が零なので
\[
 a_n=A+B(n-1)+D(n-1)(n-2)
\]
となる. 重根1に対応する多項式は,\,足し算を重ねた数列でもある.

零が$m_0$重根の場合には,\,最初の$m_0$項だけ自由でその後が零になる数列も解になる.
特性多項式を$\chi(X)=X^{m_0}Q(X)$とすると,
$Q(E)$と$E^{m_0}$は交換でき,\,そのような数列は$E^{m_0}a=0$を満たすためである.
非零根の成分にこれらを加えれば,\,初期条件の自由度は合計$k$になる.
零の冪や負の指数へ無理に代入しないことが大切である.

\subsection{非同次の式と一般化の範囲}

これまでのように項の線形結合だけで書ける式を\textbf{同次}という.
そこに外から指定された付加項がある式を\textbf{非同次}という.
右辺に既知の数列$f_n$が加わる
\[
 a_{n+k}=\sum_{j=0}^{k-1}c_j a_{n+j}+f_n
\]
でも,\,一つの解$u_n$が分かれば,\,差$b_n=a_n-u_n$は
\[
 b_{n+k}=\sum_{j=0}^{k-1}c_j b_{n+j}
\]
を満たす.
従って解は,\,\textbf{一つの特解と同次の一般解の和}になる.
不動点を引く方法は,\,特解として定数列を選んだ場合だった.

例えば
\[
 a_{n+2}-3a_{n+1}+2a_n=1
\]
では$u_n=-n$が特解である.
従って$a_n=A+B2^{n-1}-n$と書ける.
更新の異なる部分を先に一つ引き受けてもらい,\,残りを分解するのである.

なお,\,係数が$n$によって変わる場合も
$\boldsymbol v_{n+1}=C_n\boldsymbol v_n$と記録できるが,
\[
 \boldsymbol v_n=C_{n-1}\cdots C_2C_1\boldsymbol v_1
\]
となる. 毎回の行列を個別に対角化しても,\,同じ座標が次の更新に適するとは限らない.
本節の一定倍率への分解が,\,そのまま使えるわけではない.

\begin{bookpoint}
1階・高階・連立という式の違いは,\,何を状態として記録したかの違いでもある.
定数係数の線形漸化式では,\,状態を一つの行列で更新し,\,その倍率と成分間の寄与を追う.
\end{bookpoint}
